\newpage
\phantomsection
\label{prf:sqrt2-irrational}
\LRAProofFor{thm:sqrt2-irrational}

\begin{remark*}[Return]
\hyperref[thm:sqrt2-irrational]{Return to Theorem}
\end{remark*}

\begin{theorem*}[No Rational Number Has Square $2$]
There is no rational number whose square is $2$.
\end{theorem*}

\begin{proof}[Professional Standard Proof]
\LRAProofBodyStart
% MIGRATION TODO: Existing proof content preserved; detailed/professional companion body needs review.
Suppose, for contradiction, that there exists a rational number $r$ such that
\[
r^2=2.
\]
By \hyperref[rem:reduced-rational]{Remark~\ref*{rem:reduced-rational}}, we may
write $r$ in reduced form as
\[
r=\frac{p}{q},
\]
where $p,q\in\mathbb{Z}$, $\gcd(p,q)=1$, and $q>0$.

Substituting into $r^2=2$ gives
\[
\left(\frac{p}{q}\right)^2=2,
\]
so
\[
p^2=2q^2.
\]
Hence $p^2$ is even. Therefore $p$ is even, so there exists $k\in\mathbb{Z}$
such that
\[
p=2k.
\]
Substituting this into the equation above yields
\[
(2k)^2=2q^2,
\]
that is,
\[
4k^2=2q^2.
\]
Dividing by $2$, we obtain
\[
q^2=2k^2.
\]
Thus $q^2$ is even, and therefore $q$ is even.

We have shown that both $p$ and $q$ are even. But then $2$ is a common divisor
of $p$ and $q$, contradicting the assumption that $\gcd(p,q)=1$.

This contradiction proves that no rational number has square $2$.
\end{proof}

\begin{proof}[Detailed Learning Proof]
\LRAProofBodyStart
We expand the professional proof by following the same sequence of justified moves.

Suppose, for contradiction, that there exists a rational number $r$ such that
\[
r^2=2.
\]
By \hyperref[rem:reduced-rational]{Remark~\ref*{rem:reduced-rational}}, we may
write $r$ in reduced form as
\[
r=\frac{p}{q},
\]
where $p,q\in\mathbb{Z}$, $\gcd(p,q)=1$, and $q>0$.

Substituting into $r^2=2$ gives
\[
\left(\frac{p}{q}\right)^2=2,
\]
so
\[
p^2=2q^2.
\]
Hence $p^2$ is even. Therefore $p$ is even, so there exists $k\in\mathbb{Z}$
such that
\[
p=2k.
\]
Substituting this into the equation above yields
\[
(2k)^2=2q^2,
\]
that is,
\[
4k^2=2q^2.
\]
Dividing by $2$, we obtain
\[
q^2=2k^2.
\]
Thus $q^2$ is even, and therefore $q$ is even.

We have shown that both $p$ and $q$ are even. But then $2$ is a common divisor
of $p$ and $q$, contradicting the assumption that $\gcd(p,q)=1$.

This contradiction proves that no rational number has square $2$.
\end{proof}

\begin{remark*}[Proof structure]
The proof is the classical lowest-terms contradiction. Assume a rational square
root of $2$ exists and write it as $p/q$ in reduced form. The equation
\[
p^2=2q^2
\]
forces $p$ to be even. Substituting $p=2k$ then forces $q$ to be even as well.
That contradicts the assumption that $p/q$ was written in lowest terms. The
argument isolates the arithmetic obstruction behind the gap at $\sqrt{2}$.
\end{remark*}

\begin{dependencies}
\begin{itemize}
  \item \hyperref[rem:reduced-rational]{Remark~\ref*{rem:reduced-rational}} for reduced-form representatives of rational numbers
  \item Elementary parity facts for integers: if $n^2$ is even, then $n$ is even
\end{itemize}
\end{dependencies}

